\documentclass[11pt]{article}
\usepackage[T1]{fontenc}
\usepackage[utf8]{inputenc}
\usepackage{lmodern,amsmath,amssymb,amsthm,mathtools}
\usepackage[a4paper,margin=29mm]{geometry}
\usepackage{microtype,enumitem,booktabs}
\usepackage[colorlinks=true,allcolors=blue,breaklinks=true]{hyperref}
\setlength{\emergencystretch}{2em}
\newtheorem{theorem}{Theorem}[section]
\newtheorem{proposition}[theorem]{Proposition}
\newtheorem{lemma}[theorem]{Lemma}
\newtheorem{corollary}[theorem]{Corollary}
\theoremstyle{definition}\newtheorem{definition}[theorem]{Definition}
\theoremstyle{remark}\newtheorem{remark}[theorem]{Remark}
\newcommand{\OO}{\mathcal O}
\newcommand{\F}{\mathbf F}
\newcommand{\Spec}{\operatorname{Spec}}
\newcommand{\Hom}{\operatorname{Hom}}
\newcommand{\End}{\operatorname{End}}
\newcommand{\GL}{\operatorname{GL}}
\newcommand{\Jac}{\operatorname{Jac}}
\newcommand{\workbenchurl}{https://github.com/KokunoYumeto/schemes-and-moduli-workbench}
\newcommand{\recordurl}{https://github.com/KokunoYumeto/mathematics-workbenches}
\title{Grothendieck's locally quasi-finite representability criterion and Murre's announced analogue}
\author{}
\date{}
\begin{document}
\maketitle
\vspace{-2em}
\begin{abstract}
Grothendieck proposed a criterion for representing functors by separated locally quasi-finite schemes, extending the unramified criterion published by Murre in 1965. We distinguish universal properties on all local schemes from those on Artinian schemes, and specify the base of each property. Nilpotents in a henselization give a counterexample to the criterion in every positive characteristic, even with propagation relative to each complete local base. Replacing that premise by relative pro-modularity gives a criterion over every locally noetherian base. Relative propagation also makes the modularity locus open; propagation tested only over the original base does not. A second functor demonstrates this distinction and shows that effectivity for every noetherian adic completion cannot replace propagation. Neither counterexample contradicts Murre's unramified theorem. We also prove Grothendieck's finite semilocal descent theorem, including arbitrary rank in its Artinian assertions. The mathematics and text were produced and checked by GPT-6 Astra (ChatGPT Pro mode); no human has checked them.
\end{abstract}

\section{The criterion and its local universal property}

A representability criterion turns formal deformation data into a scheme. The local models in that passage must be specified carefully: a pair may represent every Artinian deformation without representing all local test schemes. Grothendieck's Montpellier notes, cote 131, sheet 13, state a criterion for a scheme ``localement quasi-fini sur $S$ et s\'epar\'e sur $S$'' \cite{Groth131}. In condition 7 of that sheet the premise is modularity on local tests. An earlier pro-modularity formulation is crossed out. Sheets 34 and 36 define the local universal property and compare it with isomorphisms of local rings.

Murre's published theorem concerns separated unramified functors \cite[Theorem 1, p.~246]{Murre}. After Corollary 2 he announces a statement with ``unramified'' replaced by ``locally quasi-finite'' \cite[p.~247]{Murre}. That announcement gives neither a complete list of modified axioms nor a proof of the locally quasi-finite statement. The distinction in the manuscript is therefore a question about its proposed extension, not a defect in Murre's published unramified criterion.

The role of henselian tests also appears in Feld's recent AI-assisted study of Morel's purity conjecture. Its Theorem~B identifies injectivity of restriction for certain Nisnevich sheaves with Gersten injectivity at henselizations \cite[Theorem~B, p.~2]{FeldPurity}. That comparison concerns sheaf stalks. Here henselization instead occurs inside the definition of a functor on all rings, and descent for that functor is a substantive part of the counterexample. None of the $K$-theoretic claims of that preprint is an input below.

We give the definitions before stating the counterexamples. Throughout, $S$ is locally noetherian, and $F$ is a contravariant set-valued functor on all $S$-schemes. Write $F(R)$ for $F(\Spec R)$ when the base map is understood.

\begin{definition}\label{def:modular}
A pair $(T,\eta\in F(T))$ is \emph{locally modular at $t\in T$} if its local scheme $\Spec\OO_{T,t}$ has the following universal property. For every local $S$-scheme $(T',t')$, every embedding $\kappa(t)\to\kappa(t')$ over the residue field of the base point, and every $\xi\in F(T')$ inducing the corresponding image of $\eta(t)$, there is a unique local map $T'\to\Spec\OO_{T,t}$ inducing that embedding and pulling $\eta$ back to $\xi$.
In addition, $\kappa(t)$ is a minimal field of definition of $\eta(t)$: if its image agrees with an element defined over another field $L$ in a common overfield, the image of $\kappa(t)$ is contained in the image of $L$, and that element is the induced one. Restricting the universal property to local Artinian test schemes gives \emph{pro-modularity}.
\end{definition}

For a represented functor $F=h_X$, local modularity says exactly that $\OO_{X,x}\to\OO_{T,t}$ is an isomorphism, with the specified residue field. This is the interpretation stated on sheet 36 of the manuscript. It includes tests whose local rings are not Artinian.

Here is a formulation of the eight conditions that specifies the base in each local universal property. In condition~\ref{ax:propagate} we first use the original base $S$; in condition~\ref{ax:generic} we use the displayed new base.

\begin{enumerate}[label=(\arabic*),ref=\arabic*,leftmargin=*]
\item\label{ax:zariski} $F$ is a Zariski sheaf.
\item\label{ax:fpqc} $F$ is a sheaf for faithfully flat quasi-compact covers.
\item\label{ax:limits} $F$ commutes with filtered colimits of rings.
\item\label{ax:complete} For every complete noetherian local $S$-algebra $(A,\mathfrak m)$, the map $F(A)\to\varprojlim_n F(A/\mathfrak m^{n+1})$ is bijective.
\item\label{ax:artin} Over every local Artinian $S$-algebra, the local deformation functors are universally prorepresented by finite local algebras, including the minimal fields of definition and compatibility with residue-field extensions. Equivalently, there is a disjoint union of spectra of finite local algebras whose functor agrees on all local Artinian tests, universally for residue extensions.
\item\label{ax:valuation} For every valuation-ring test $V$, the map $F(V)\to F(\operatorname{Frac}V)$ is injective.
\item\label{ax:propagate} Let $A$ be a complete noetherian local $S$-algebra, $A\to B$ a local homomorphism with $B$ finite over $A$, and $\xi\in F(B)$. If this pair is locally modular for $F$ on $S$-schemes at the closed point of $\Spec B$, it is locally modular for that same functor at all points.
\item\label{ax:generic} Let $S'$ be irreducible noetherian over $S$, let $Z$ be locally of finite type over $S'$, and let $\eta\in F(Z)$ be locally modular for the restricted functor $F_{S'}$ on $S'$-schemes at a point over the generic point of $S'$. Relative local modularity then holds on a nonempty open of the corresponding irreducible component near that point.
\end{enumerate}

The finite-type qualification on $Z$ in condition~\ref{ax:generic} is explicit here. It is absent from the displayed wording on sheet 13. Without it even necessity fails for a represented functor, as Proposition~\ref{prop:finite-type} shows. This qualification does not repair the sufficiency problem addressed below.

We also need a relative replacement for condition~\ref{ax:propagate}. Write $(7_{\mathrm{rel}})$ for the assertion with the same $A,B,\xi$, but with both the premise and conclusion formulated for the restricted functor $F_A$ on $A$-schemes. Its alternative conclusion is that $h_{\Spec B}\to F_A$ is representable and \'etale. Sheet 13 writes $F(B)$ in condition 7, whereas condition 8 explicitly writes the restricted functor over $S'$. We distinguish the two formulations without assuming an implicit change of base in the first. The distinction is substantive, as Theorem~\ref{thm:adic} shows. A local universal model after completing the base need not be a local universal model over $S$.

\begin{theorem}\label{thm:counterexample}
For every prime $p$, over $S=\Spec\F_p[t]$, the affine formula
\[
 P(R)=\{x\in R^h_{(t)}:x^p=0\}
\]
extends by Zariski gluing to a functor satisfying conditions \ref{ax:zariski}--\ref{ax:generic} and $(7_{\mathrm{rel}})$. Here $R^h_{(t)}$ is the henselization of the pair $(R,tR)$. Nevertheless $P$ is not represented by any separated $S$-scheme. Its closed fibre at $t=0$ is $\alpha_p$, so it is not an unramified functor.
\end{theorem}

The reason the propagation condition fails to force representability is concrete. Over $\F_p[[t]]$ the Artinian model is $\alpha_p$. A local test outside the Artinian category acquires an Artin--Schreier root in its henselization, and multiplying that root by a square-zero element gives an extra nilpotent section. Hence this Artinian model is not a local universal model for the restricted functor. For modularity over $S$ itself, the zero section on the local base gives an even more direct obstruction: a finite complete local model could exist only over the generic point of $S$. The substantial descent issue is that a henselization of a tensor product need not equal the tensor product of the henselizations. Section~\ref{sec:descent} proves the needed arbitrary fpqc descent directly.

\begin{theorem}\label{thm:adic}
For every prime $p$, define a subfunctor of $\alpha_p$ on $\F_p[t]$-algebras by
\[
 H(R)=\{x\in R:x^p=0,\ \text{there is }a\in R\text{ with }x=tax\}
\]
This is an fpqc sheaf. For every noetherian ring $R$ complete and separated for an arbitrary ideal $I$, it satisfies
\[
 H(R)\simeq\varprojlim_n H(R/I^n).
\]
Over $S_0=\Spec\F_p[t]$, it satisfies conditions \ref{ax:zariski}--\ref{ax:generic}, but the modularity locus of $(S_0,0)$ is the singleton $\{(t)\}$ and is not open. After restriction to $\Spec A_0$, where $A_0=\F_p[[t]]$, it satisfies conditions \ref{ax:zariski}--\ref{ax:valuation} and \ref{ax:generic}, but fails condition~\ref{ax:propagate}: the pair $(\Spec A_0,0)$ is locally modular at the closed point and not at the generic point. Consequently the functor over $S_0$ also fails $(7_{\mathrm{rel}})$.
\end{theorem}

\begin{theorem}\label{thm:openness}
Assume conditions \ref{ax:zariski}--\ref{ax:valuation}, $(7_{\mathrm{rel}})$, and \ref{ax:generic}. For every locally finite-type pair $(T,\eta)$, the local modularity locus $M(T,\eta)$ for $F$ on $S$-schemes is open. No representation of $F$ over the local rings of $S$ is assumed. The \'etale alternatives in the conclusions of $(7_{\mathrm{rel}})$ and condition~\ref{ax:generic} also suffice.
\end{theorem}

The positive criterion uses relative pro-modularity. Call condition $(7_{\mathrm{pro}})$ the following assertion: if $A$ is a complete noetherian local $S$-algebra, $B$ is finite local over $A$, and $\xi\in F(B)$ represents the local Artinian deformation functor of $F_A$ at its closed point, including the specified residue fields, then $\Spec B\to F_A$ is locally modular at every point. The alternative conclusion that this map is representable and \'etale is sufficient as well. This condition implies $(7_{\mathrm{rel}})$: a pair locally modular for $F_A$ has the required universal property on Artinian $A$-tests.

\begin{theorem}\label{thm:criterion}
Over every locally noetherian scheme $S$, conditions \ref{ax:zariski}--\ref{ax:valuation}, $(7_{\mathrm{pro}})$, and \ref{ax:generic} are equivalent to representability by a separated locally quasi-finite $S$-scheme. No excellence assumption on $S$ is imposed.
\end{theorem}

Thus even effectivity for all noetherian adic completions cannot replace a propagation axiom, and the base of that axiom matters. Theorems~\ref{thm:counterexample} and \ref{thm:adic} concern functors of points, including nonreduced tests. Their field-valued points alone carry too little information to detect either obstruction.


The descent mechanism also has a finite algebra counterpart. Grothendieck's cote 128, sheet 25, asserts a ``morphisme de descente effective'' for flat modules from cohomological vanishing on Artinian quotients \cite{Groth128}. We prove the statement in every rank in its Artinian assertions, while distinguishing finite rank over a general semilocal base.

\begin{theorem}\label{thm:finite-descent}
Let $A$ be a semilocal noetherian ring, $\mathfrak m=\operatorname{rad}A$, and $A'$ a finite $A$-algebra. Put $A_n=A/\mathfrak m^{n+1}$, $A'_n=A'\otimes_AA_n$, and
\[
 B=H^0(A'/A),\qquad B^{(n)}=H^0(A'_n/A_n),
\]
where $H^0$ is the equalizer of the two maps to the tensor square. For any fixed integer $r\geq1$, the following three conditions are equivalent:
\begin{enumerate}[label=(\roman*),leftmargin=*]
\item $H^1(A'_n/A_n,\mathbf G_a)=0$ for every $n$;
\item $H^1(A'_n/A_n,\GL_r)=\{1\}$ for every $n$;
\item $\Spec A'_n\to\Spec B^{(n)}$ is effective descent for all flat quasi-coherent modules for every $n$.
\end{enumerate}
Here $H^1$ denotes Amitsur cohomology, or its matrix-cocycle pointed set. If these conditions hold, then they hold for every finite $r$; moreover
\[
 H^1(A'/A,\mathbf G_a)=0,\qquad
 H^1(A'/A,\GL_r)=\{1\}\quad(r\geq1),
\]
and $\Spec A'\to\Spec B$ is effective descent for finite flat modules. If $A$ is Artinian, this map and its base change by every flat $B$-scheme are effective descent for all flat quasi-coherent modules, with no bound on rank.
\end{theorem}

The proof removes a matrix cocycle one nilpotent order at a time. Additive primitives can be normalized to vanish modulo the radical, so the correcting matrices are invertible even in arbitrary rank. Conversely, scalar matrix cocycles recover additive vanishing, including from rank one. Artin--Rees turns a sufficiently accurate finite-rank basis into an exact basis over the semilocal ring. For flat base changes, descent of the scalar action and a nilpotent flatness criterion finish the argument.

\begin{proposition}\label{prop:finite-type}
If condition~\ref{ax:generic} allows arbitrary $Z$, it need not hold even for the identity functor $h_S$.
\end{proposition}
\begin{proof}
Let $k$ be algebraically closed, $A=k[t]$, and $M=\bigoplus_{a\in k}A/(t-a)$. Give $B=A\oplus M$ the square-zero multiplication, and take $Z=\Spec B$ over $S=\Spec A$. It has the same underlying space as $S$. At the generic point $M$ vanishes, and the local-ring map is the identity of $k(t)$. At the point $t=a$ the map is $A_{(t-a)}\to A_{(t-a)}\oplus k$, which is not an isomorphism. Thus its modularity locus is precisely the generic point and contains no nonempty open.

For a locally finite-type test over a noetherian base, an isomorphism of local rings of a morphism locally of finite presentation spreads to an isomorphism on neighborhoods: it first makes the morphism \'etale there, then finite presentation spreads the inverse map, whose image is open. This proves the corrected necessity used above.
\end{proof}

\section{A local algebra descent lemma}\label{sec:localalgebra}

We prove the algebraic fact that permits descent of a nilpotent section even when its equality on overlaps is known only after henselization. No noetherian hypothesis is used in this section.

\begin{lemma}\label{lem:finite-module}
Let $A\to B$ be flat and $b\in B$ integral over $A$. Suppose there is $e\in B\otimes_A B$ with
\[
 (b\otimes1-1\otimes b)e=0,\qquad \mu(e)=1,
\]
where $\mu$ is multiplication. There is a finitely presented $A$-module $M$ with an endomorphism $T$ such that the $B$-module of rank one on which $T$ acts by $b$ is a direct summand of $M\otimes_A B$, compatibly with $T$.
\end{lemma}
\begin{proof}
Choose a monic equation $f(b)=0$ and put $C=A[T]/(f)$. Write $e=\sum_{i=1}^r u_i\otimes v_i$. Map $C^r\to B$ by sending its basis vectors to the $u_i$, and call the kernel $K$. The vector $e_0=\sum_i e_i\otimes v_i$ maps to $e$. Since $B$ is $A$-flat, the error $(T-b)e_0$, which maps to zero, belongs to $K\otimes_A B$. It is a finite sum of tensors $k_j\otimes w_j$. Let $K_0$ be the $C$-module generated by the finitely many $k_j$ and put $M=C^r/K_0$.

The algebra $C$ is finite free over $A$, so $M$ is finitely presented as an $A$-module. The image $e_M$ of $e_0$ satisfies $(T-b)e_M=0$. The induced map $M\otimes_A B\to B$ sends $e_M$ to $\mu(e)=1$. Therefore $z\mapsto e_Mz$ is a section respecting $T$, as claimed.
\end{proof}

\begin{lemma}\label{lem:indecomposable}
Let $(A,\mathfrak m,k)$ be strictly henselian local, and let $M$ be a finitely presented $A$-module with an endomorphism $T$. It has a finite $T$-stable direct-sum decomposition $M=\bigoplus_iM_i$ such that $\End_{A[T]}(M_i)$ is local. For every flat local map $A\to(B,\mathfrak n,K)$, the ring $\End_{B[T]}(M_i\otimes_A B)$ is also local.
\end{lemma}
\begin{proof}
For a finite $A$-module $N$ with $T$ action, let $E=\End_{A[T]}(N)$ and let $J$ be the kernel of its action on $N/\mathfrak mN$. Its image $E/J$ is finite dimensional over $k$. Moreover $J\subset\Jac(E)$. For $h\in J$ and $g\in E$, Nakayama's lemma makes $1-gh$ surjective on $N$. A surjective endomorphism $u$ of a finite module over a commutative ring is invertible: if generators satisfy $x_i=u(\sum_j a_{ij}x_j)$, the determinant identity gives an equation $1+u q(u)=0$ on $N$. Thus $-q(u)$ is its inverse. This proves the radical assertion without a noetherian assumption.

Idempotents of $E/J$ lift to $E$. Choose a minimal generating set of $N$ and a surjection $A^r\to N$. Lift an endomorphism $h$ whose reduction is idempotent to a matrix $H$ on $A^r$. Its characteristic polynomial reduces to $X^a(X-1)^b$. Henselian factorization lifts these relatively prime monic factors. The Chinese remainder theorem supplies a polynomial $q$ equal to zero on the first factor and one on the second. Cayley--Hamilton makes $q(H)$ idempotent. It preserves the kernel of $A^r\to N$, so induces the required idempotent $q(h)$. Splitting a finite primitive-idempotent decomposition of $E/J$, successively in complementary summands, gives $N=\bigoplus_iN_i$. Each corner has local quotient and radical kernel, and is therefore local.

For the assertion after base change, finite presentation and flatness give
\[
 \End_{A[T]}(M_i)\otimes_A B
   \simeq\End_{B[T]}(M_i\otimes_A B).
\]
One obtains this by applying $\Hom$ to a finite presentation, then taking the kernel of the commutator with $T$. Let $R_i$ be the image of the left-hand endomorphism ring on $M_i/\mathfrak mM_i$. It is a finite-dimensional local $k$-algebra. Its residue division algebra is a finite purely inseparable field extension $D_i$ of $k$: its center has that form because $k$ is separably closed, and a central division algebra over a separably closed field splits and is the field itself.

The image after base change on the residue module is $R_i\otimes_k K$. Its radical contains the nilpotent ideal $\Jac(R_i)\otimes_k K$, with quotient $D_i\otimes_k K$. A finite purely inseparable field extension tensored with any field is local Artinian, since each generator has only one possible image in an algebraic closure. Thus the residue image is local. Its kernel is in the Jacobson radical by the finite-module argument already proved over $B$. The whole endomorphism ring is local.
\end{proof}

\begin{proposition}\label{prop:integral-descent}
If $A$ is strictly henselian local, $A\to B$ is flat local, and an integral element $b\in B$ admits a tensor $e$ as in Lemma~\ref{lem:finite-module}, then $b$ comes from a unique element of $A$.
\end{proposition}
\begin{proof}
Place the rank-one module $B_b$ as a summand of $M_B$ using Lemma~\ref{lem:finite-module}, and decompose $M$ by Lemma~\ref{lem:indecomposable}. If $s_i:B_b\to(M_i)_B$ and $r_i:(M_i)_B\to B_b$ are the components of the inclusion and retraction, then $\sum_i r_is_i=1$ in $B$. Since $B$ is local, one summand is a unit. After rescaling, $B_b$ is a nonzero direct summand of that $(M_i)_B$. Its local endomorphism ring has no nontrivial idempotents, so $(M_i)_B\simeq B_b$.

The flat local map is faithfully flat. Reduction to the residue field shows that $M_i$ has one generator. A surjection $A\to M_i$ becomes an isomorphism after tensoring with $B$, because it is a surjection between free rank-one $B$-modules. Faithful flatness kills its kernel, so $M_i\simeq A$. The endomorphism $T$ is multiplication by an element $a\in A$, and the isomorphism respecting $T$ gives $b=a$ in $B$. Injectivity of $A\to B$ proves uniqueness.
\end{proof}

\section{Fpqc descent for nilpotents in a henselization}\label{sec:descent}

Write $R^h$ for the henselization of $(R,tR)$. We first record the elementary properties needed to apply the preceding algebra.

\begin{lemma}\label{lem:hensel}
An element $z\in R$ maps to zero in $R^h$ if and only if $(1+tc)z=0$ for some $c\in R$. If $t\in\Jac(R)$ and $R\to D$ is flat and faithfully flat modulo $t$, then it is faithfully flat. A henselian local ring is henselian along every ideal contained in its maximal ideal.
\end{lemma}
\begin{proof}
In the second assertion, the image of $\Spec D$ contains every maximal ideal of $R$, and the image of a flat map is stable under generization. It therefore contains every prime. For the first, localize $R$ at $1+tR$. In this localization $t$ belongs to the Jacobson radical. The map to $R^h$ is flat and is an isomorphism modulo $t$, so is faithfully flat by the preceding argument. Its kernel is zero. The kernel before localization is precisely the asserted annihilator condition.

For the last assertion, every finite algebra over a henselian local ring is a finite product of local rings. Reducing each factor modulo an ideal in the maximal ideal preserves its two idempotents. Hence reduction gives a bijection on idempotents in every finite algebra, which is the henselian-pair criterion \cite[Tags 04GG, 09XI]{Stacks}.
\end{proof}

The functor $P$ commutes with filtered ring colimits. Henselization of pairs does so by its universal property: every monic factorization problem and its coprimality witnesses occur at a finite stage, as does the invertibility of each $1+tb$. Taking the kernel of $x\mapsto x^p$ also commutes with filtered colimits.

It is an \'etale sheaf. Indeed, after replacing the base ring by its henselization, the henselizations of an \'etale cover and of its tensor powers are unchanged by the universal property. Apply de Jong--Olander's published henselian descent theorem to the scheme $\alpha_p$ \cite[Theorem 2, p.~2]{dJO}. Its assumption is a weakly \'etale faithfully flat map, and an \'etale cover is such a map. In particular finite affine Zariski covers give Zariski locality, and gluing defines $P$ on all schemes.

For any faithfully flat $R\to S$, the map $R^h\to S^h$ is flat, and modulo $t$ it is faithfully flat. Since $t\in\Jac(R^h)$, Lemma~\ref{lem:hensel} makes it faithfully flat. In particular $P(R)\to P(S)$ is injective.

\begin{theorem}\label{thm:fpqc}
For every faithfully flat map of $\F_p[t]$-algebras $R\to S$, the diagram
\[
 P(R)\longrightarrow P(S)\rightrightarrows P(S\otimes_R S)
\]
is an equalizer. Consequently $P$ is an fpqc sheaf on schemes.
\end{theorem}
\begin{proof}
First let $A$ be strictly henselian local and $A\to R$ faithfully flat, and suppose $b\in P(R)$ has equal images on the overlap. If $t$ is a unit in $A$, all the henselizations are zero and there is nothing to prove. Otherwise $(A,tA)$ is henselian. Put $B=R^h$. Then $A\to B$ is faithfully flat by Lemma~\ref{lem:hensel}, and $b^p=0$.

The overlap equality maps to $(B\otimes_A B)^h$. The kernel description supplies $c\in B\otimes_A B$ with
\[
 (1+tc)(b\otimes1-1\otimes b)=0.
\]
Its multiplication image $u=1+t\mu(c)$ is a unit of $B$. Thus $e=(u^{-1}\otimes1)(1+tc)$ has multiplication image one and annihilates the difference. Choose a prime $\mathfrak q$ of $B$ above the maximal ideal of $A$. The map $A\to B_{\mathfrak q}$ is flat local. Proposition~\ref{prop:integral-descent}, applied to the image of $e$, gives $a\in A$ with $b=a$ in $B_{\mathfrak q}$.

This equality holds on all of $B$. Indeed, put $D=(B\otimes_A B_{\mathfrak q})^h$. The map $B\to D$ is flat and faithfully flat modulo $t$: modulo $t$ it is the base change of the faithfully flat local map $A/tA\to B_{\mathfrak q}/tB_{\mathfrak q}$. As $t\in\Jac(B)$, it is faithfully flat. The cocycle equality in $D$ identifies $b$ from the first factor with $a$ from the second. Injectivity gives $b=a$ in $B$, and faithful flatness of $A\to B$ gives $a^p=0$. Thus the section descends over a strictly henselian local base.

For general $R\to S$, base change at every geometric point of $\Spec R$ to the strict henselization $R_{\mathfrak p}^{\mathrm{sh}}$. The result just proved gives a section there whose pullback is the given section of $P(S)$. Finitarity descends that section to a pointed finitely presented \'etale neighborhood, and, after increasing the stage, descends its equality with the section on $S$ as well. These neighborhoods cover $\Spec R$. Their sections agree on overlaps by the fpqc injectivity already proved, because they agree after base change by $S$. The \'etale sheaf property glues them to a unique element of $P(R)$. This proves the affine equalizer. A quasi-compact flat surjection over an affine scheme can be refined by finitely many affine opens in its source, whose disjoint union is affine. The affine assertion and Zariski gluing therefore prove the stated fpqc sheaf assertion.
\end{proof}

\section{The counterexample to representability}

We now check every remaining condition and prove Theorem~\ref{thm:counterexample}. On a complete noetherian local ring $A$, either $t$ is a unit, so $P(A)$ is the one-element set, or $t$ belongs to the maximal ideal and $(A,tA)$ is henselian, so $P(A)=\alpha_p(A)$. The same alternatives hold in the Artinian quotients. Completeness of $A$ proves condition~\ref{ax:complete}.

Over an Artinian local base $C$, $t$ is either invertible or nilpotent. In the former case the entire restricted functor is terminal. In the latter case every $C$-algebra is henselian along the nilpotent ideal generated by $t$, so the entire restricted functor is $\alpha_{p,C}$. These are finite local representing schemes, universally under residue-field extensions. This proves condition~\ref{ax:artin} with more than just its Artinian-test assertion. Henselization of a reduced ring is reduced, since it is a filtered colimit of \'etale neighborhoods. Thus $P$ has only its zero section on valuation rings and fields; condition~\ref{ax:valuation} follows.

\begin{lemma}\label{lem:no-model}
Let $(A,\mathfrak m)$ be a complete noetherian local $\F_p[t]$-algebra with $t\in\mathfrak m$ not nilpotent. There is no finite local $A$-algebra $B$ carrying a section that is locally modular at its closed point for the restricted functor $P_A$ on $A$-schemes.
\end{lemma}
\begin{proof}
Suppose $(B,\xi)$ were such a pair. There is only one point over each field, defined over the base residue field, so the minimal-field condition forces $\kappa(B)=A/\mathfrak m$. On every Artinian local $A$-test centered there, $P$ is $\alpha_p$. The two local universal properties identify the complete local rings and their elements:
\[
 (B,\xi)\simeq(A[x]/(x^p),x).
\]
For detail, apply the two properties to the Artinian quotients of the rings, obtain compatible maps in both directions, and pass to the limits. Their composites are identities on each quotient and therefore on the complete rings.

Take $R_0=A[u]_{(\mathfrak m,u)}$ and $R=R_0[\epsilon]/(\epsilon^2)$. In $R_0^h$ there is a unique root
\[
 f^p-f=tu,\qquad f\equiv0\pmod t.
\]
It is not in $R_0$. Otherwise reduce modulo a minimal prime of $A$ avoiding $t$, which exists because $t$ is not nilpotent. In the rational function field $\operatorname{Frac}(A/\mathfrak q)(u)$ this would give an Artin--Schreier root of $tu$. At infinity the right side has a pole of order one, while every pole of $g^p-g$ has order divisible by $p$. This is impossible.

Henselization commutes with integral base change, so $R^h=R_0^h[\epsilon]/(\epsilon^2)$. The element $\epsilon f$ belongs to $P(R)$ and reduces to the unique field-valued point. It does not come from $\alpha_p(R)$: comparing its $\epsilon$ coefficient would place $f$ in $R_0$. The comparison is valid because $R_0\to R_0^h$ is faithfully flat, hence injective. This contradicts local modularity of the identified pair $A[x]/(x^p)$ on the local test $R$.
\end{proof}

We check condition~\ref{ax:propagate} with modularity over its stated base $S$. Suppose $(B,\xi)$ is locally modular at its closed point $b$, and write $s$ for its image in $S$. Since there is only one field-valued section, the minimal-field condition gives $\kappa(b)=\kappa(s)$. The ring $B$ is complete: it is finite over the complete local ring $A$, and the $\mathfrak m_A B$-adic and maximal-ideal topologies agree. Put $R=\mathcal O_{S,s}$. Its zero section is an eligible local test, so modularity gives a local $S$-map $B\to R$. The structural map to $B$ factors through $R$, and the composite $R\to B\to R$ is the identity. Thus $R$ is a quotient of $B$ and is complete.

If $s$ were closed, $R=\F_p[t]_{(g)}$ for an irreducible polynomial $g$. This ring is countable, whereas its completion, a power-series ring over its finite residue field, is uncountable. Thus $s$ is generic and $R=\F_p(t)$. Here $P$ is terminal. Modularity and uniqueness on the local test $B$ make $R\to B$ and $B\to R$ inverse. Hence $B=\F_p(t)$ has only one point, and condition~\ref{ax:propagate} follows.

For $(7_{\mathrm{rel}})$, let $A$ be complete noetherian local. If $t$ is a unit, $P_A$ is terminal; if $t$ is nilpotent, $P_A=\alpha_{p,A}$. In either case its representing scheme is finite local over $A$. The local universal property at the closed point identifies any proposed finite local model $B$, together with its element, with that entire representing scheme. Thus the pair is locally modular everywhere. In the remaining case, $t$ is a nonnilpotent element of the maximal ideal, and Lemma~\ref{lem:no-model} shows that the premise cannot hold. This proves the relative condition in all cases.

For condition~\ref{ax:generic}, let $S'$ be irreducible noetherian. If $t$ vanishes at its generic point, it is locally nilpotent, and the functor is $\alpha_p$. Otherwise $D(t)$ is a dense open containing the generic point and the functor is terminal there. Near the point in question it therefore has a separated finite-type representation. For a locally finite-type test, a local-ring isomorphism spreads to an open isomorphism as proved in Proposition~\ref{prop:finite-type}. This proves condition~\ref{ax:generic}.

Finally put $R=\F_p[t,\epsilon]/(\epsilon^p)$. The sections $0$ and $\epsilon$ of $P(R)$ are distinct, as their reductions modulo $t$ are distinct. They become equal on $D(t)$, where $P$ is terminal. The inclusion $D(t)\subset\Spec R$ is schematically dense, since $R$ is free over $\F_p[t]$ and multiplication by $t$ is injective. If a separated $S$-scheme represented $P$, the two sections would define maps whose equalizer is a closed subscheme containing this schematically dense open. That equalizer is all of $\Spec R$, a contradiction. This excludes every separated representing scheme and completes Theorem~\ref{thm:counterexample}.

\section{Adic effectivity and the base of modularity}

We prove Theorem~\ref{thm:adic} by testing separately the equations and the local universal property. The formula for $H$ is functorial and defines a subgroup of $\alpha_p$. If $(1-ta)x=(1-tb)y=0$, then
\[
 (1-t(a+b-tab))(x+y)=(1-ta)(1-tb)(x+y)=0.
\]
Additive inverses satisfy the same equation.

For a faithfully flat map $R\to R'$, a compatible element of $H(R')$ first descends uniquely to $x\in\alpha_p(R)$. Put $M=Rx$. Flatness identifies its base change with $R'x$, and faithful flatness detects whether the class of $x$ in $M/tM$ is zero. The equation over $R'$ makes that class zero; hence $x\in tRx$ over $R$. This proves arbitrary affine fpqc descent, and Zariski gluing extends the functor to schemes. The equations $x^p=0$ and $x=tax$ use finitely many elements and operations, so both sections and their witnesses descend through filtered colimits. Equality is equality of ring elements and also descends.

Let $R$ be noetherian and complete and separated for $I$. A compatible system $x_n\in H(R/I^n)$ gives $x\in R$ with $x^p=0$. Choose a witness at each level and lift it to $a_n\in R$. No compatibility among these witnesses is required. They give
\[
 x\in (tx)+I^n\quad\text{for every }n.
\]
Completeness puts $I$ in $\Jac(R)$ by the convergent geometric series for $1-i$. Krull intersection applied to the finite module $R/(tx)$ yields $\bigcap_n((tx)+I^n)=(tx)$. Hence $x\in tRx$. Separatedness gives injectivity, proving the asserted adic effectivity for every ideal of definition.

Over an Artinian local base, if $t$ is a unit then $H=\alpha_p$, and if $t$ is nilpotent then $1-ta$ is always a unit and $H$ is terminal. This verifies finite universal Artinian representability. On a valuation ring the only section is zero. For an irreducible noetherian test base $S'$, either $t$ is locally nilpotent or it is invertible on a dense open. Thus over a dense open $U\subset S'$ the functor is represented by $U$ or $\alpha_{p,U}$, and the same finite-type spreading argument verifies condition~\ref{ax:generic}.

We next verify condition~\ref{ax:propagate} over the original polynomial base $S_0$. Suppose $B$ is finite local over a complete noetherian local algebra $A$, and $(B,\xi)$ is locally modular for $H$ on $S_0$-schemes at its closed point $b$, with image $s\in S_0$. Every field-valued object is zero, so minimality gives $\kappa(b)=\kappa(s)$. The zero object of $R=\OO_{S_0,s}$ is an eligible local test. Its universal map gives a local homomorphism $B\to R$ splitting the structural map $R\to B$. The ring $B$ is complete, since it is finite local over $A$, and every quotient of a complete noetherian local ring is complete. Thus $R$ would be complete. As in the preceding section, countability of $\F_p[t]_{(g)}$ and uncountability of its completion exclude every closed $s$.

It follows that $s$ is generic. Put $K=\F_p(t)$. All local tests centered there factor through $K$, and $H_K=\alpha_{p,K}$. The two local universal maps therefore identify $(B,\xi)$ with $(K[z]/(z^p),z)$. This has only one point, so condition~\ref{ax:propagate} holds. This verifies its all-points modularity conclusion; no \'etale conclusion for a morphism over $S_0$ is needed here.

Nevertheless the modularity locus of $(S_0,0)$ is not open. At $(t)$, every eligible local test has $t$ in its maximal ideal, so $1-ta$ is a unit and $H$ has only its zero object. The structural map to $\Spec\OO_{S_0,(t)}$ is the unique local $S_0$-map, and $\F_p$ is the minimal field of definition. At any point $s\in D(t)$ the local test $\kappa(s)[\epsilon]/(\epsilon^p)$ instead has the eligible nonzero object $\epsilon$. A pullback of the zero section cannot give that object. Hence
\[
 M(S_0,0)=\{(t)\}.
\]

Finally restrict the base to $\Spec A_0$, where $A_0=\F_p[[t]]$. The preceding sheaf, completion, Artinian and generic arguments remain valid. Take $A=B=A_0$ and $\xi=0$. Every local $A_0$-test centered at the closed point again has only the zero object, and exactly one structural map to $\Spec A_0$, with minimal residue field $\F_p$. Thus the pair is locally modular at the closed point. At the generic point, put $K=\F_p((t))$ and use the nonzero object $\epsilon\in H(K[\epsilon]/(\epsilon^p))$. It cannot be a pullback of zero. This disproves propagation over $A_0$ and also disproves $(7_{\mathrm{rel}})$ for the original functor over $S_0$. Theorem~\ref{thm:adic} follows.

\section{Artinian comparison and \'etale charts}\label{sec:comparison}

We now prove the locally quasi-finite criterion with the pro-modular premise. The first step explains exactly how Artinian tests determine a natural transformation once descent and effectivity are available.

\begin{lemma}\label{lem:comparison}
Let $G,F$ satisfy conditions \ref{ax:zariski}--\ref{ax:complete}, and let $G\to F$ be a natural transformation bijective on every local Artinian $S$-algebra. Then it is an isomorphism on all $S$-schemes.
\end{lemma}
\begin{proof}
First prove injectivity. Over a noetherian ring, two sections with equal image agree on every Artinian quotient at every maximal ideal. Condition~\ref{ax:complete} makes them agree on the corresponding complete local rings; fpqc injectivity then makes them agree on the local rings. Finitarity spreads those equalities to neighborhoods, and Zariski descent proves global equality. For an arbitrary algebra over a noetherian affine base, a pair of sections and equality of their images descend to a finite-type algebra by finitarity. This proves injectivity on arbitrary rings, including the nonnoetherian overlap rings that occur next.

For surjectivity over a noetherian local $A$, an element of $F(A)$ has unique compatible preimages over $A/\mathfrak m^n$. They give an element of $G(\widehat A)$. Its two images over $\widehat A\otimes_A\widehat A$ agree by the injectivity just proved. Fpqc descent gives an element of $G(A)$ with the prescribed image. Localizing and spreading extends the conclusion to noetherian rings, and finite-type approximation then extends it to all rings. Zariski gluing finishes the proof.
\end{proof}

A scheme locally of finite presentation satisfies the effectivity condition used here: a compatible sequence of maps from the Artinian quotients of a complete local ring lands in a fixed affine neighborhood of their common closed point. Taking the inverse limit of the coordinate homomorphisms gives the map from the complete local ring. In particular no properness hypothesis is needed for this comparison.

Condition~\ref{ax:artin}, combined with Lemma~\ref{lem:comparison}, has an important consequence: over a local Artinian base, $F$ is represented on \emph{all} test schemes by a disjoint union of finite local schemes. The finite representing algebras in condition~\ref{ax:artin} are themselves Artinian and carry their universal elements, so they give the required natural transformation before the comparison lemma is applied.

\begin{lemma}[Artinian form of the \'etale criterion]\label{lem:etale}
Over a locally noetherian base, conditions \ref{ax:zariski}--\ref{ax:complete} and \ref{ax:valuation}, together with unique lifting through nilpotent ideals on local Artinian tests, imply representability by a separated \'etale scheme.
\end{lemma}
\begin{proof}
Murre's Corollary 2 \cite[p.~247]{Murre} proves this with unique lifting on all nilpotent tests. We verify the reduction, beginning with uniqueness. For a noetherian local pair $(A,I)$, equality modulo the nilpotent ideal $I$ implies equality modulo every $\mathfrak m^n$: apply Artinian uniqueness to $A/\mathfrak m^n\to A/(I+\mathfrak m^n)$. Condition~\ref{ax:complete} and faithful flatness of completion give equality over $A$. Localization gives uniqueness for noetherian rings.

For an arbitrary algebra $A$ over a noetherian affine part of the base, write
\[
 A/I=\varinjlim_{\substack{J\subseteq I\\J\text{ finitely generated}}} A/J.
\]
An equality of two sections over $A/I$ already holds over one $A/J$ by finitarity. Choose generators of $J$ and an exponent $n$ with $I^n=0$. The two sections, their equality modulo $J$, these generators and the finitely many equations making all degree-$n$ monomials in them zero descend to a finite-type algebra. Its corresponding ideal is nilpotent. Uniqueness there proves uniqueness over $A$. Thus uniqueness holds also on the possibly nonnoetherian completion overlaps.

For existence over a noetherian local $A$, restrict a section over $A/I$ to $A/(I+\mathfrak m^n)$ and lift it uniquely to $A/\mathfrak m^n$. The lifts are compatible and give a section over $\widehat A$. Its reduction agrees with the original section over $\widehat A/I\widehat A$, by effectivity for that complete noetherian quotient. Its two pullbacks to $\widehat A\otimes_A\widehat A$ have equal reductions modulo $I$; the uniqueness already proved makes them equal. Fpqc descent gives the lift over $A$. Localization, spreading and Zariski gluing handle noetherian rings. Finally, a section over an arbitrary $A/I$ comes from one $A/J$ in the displayed limit. Descend that section, generators of $J$ and their nilpotence equations to a finite-type stage and apply the noetherian existence assertion. Pullback gives the required lift over $A$. Murre's corollary now applies; its remaining hypotheses are exactly the stated sheaf, finitary, complete-local and valuative conditions.
\end{proof}

For a finite-type pair $(T,\eta)$, let $E(T,\eta)$ be its \emph{Artinian formally \'etale locus}: at a point of $E$, every diagram of a nilpotent closed immersion of local Artinian schemes centered at that point, a map into $T$ on the closed subscheme, and an element of $F$ extending its image has a unique lifting map to $T$. Residue-field extensions are allowed. This property is invariant under \'etale changes of source and base and can be checked after such a change: lift the residue point after a finite \'etale field extension, apply formal \'etaleness of the change, and descend the unique map. Write $M(T,\eta)$ for the local modularity locus of Definition~\ref{def:modular}.

\section{Openness of the modularity locus}\label{sec:openness}

This section proves openness with propagation relative to complete local bases. The main issue is to control infinitesimal directions transverse to a closed stratum; a statement only about its reduced generic point does not do that.

\begin{lemma}\label{lem:uniform}
Let $R$ be noetherian, $I\subset R$ an ideal, and $C$ a finite $R$-module. There is $N\geq1$ such that the following holds. If $R\to R'$ is flat, $R'$ is noetherian local and $IR'$ is contained in its maximal ideal, $D$ is finite over $R'$, and $\phi:D\to C\otimes_RR'$ is an isomorphism modulo $I^N$, then $\phi$ is an isomorphism. The same $N$ works for any direct-summand target of $C\otimes_RR'$.
\end{lemma}
\begin{proof}
Choose $0\to K\to R^d\to C\to0$. Artin--Rees gives $N$ with $K\cap I^NR^d\subset IK$ \cite[Tag 00IN]{Stacks}. Flat base change preserves the presentation and this intersection. Nakayama makes $\phi$ surjective. Lift the chosen generators to $D$; they generate modulo $I^N$ and hence generate $D$. If $K'$ is the kernel of $R'^d\to D$, then $K'\subset K\otimes_RR'$ and the isomorphism modulo $I^N$ gives
\[
 K\otimes_RR'\subset K'+I^NR'^d.
\]
Intersecting with the kernel and using Artin--Rees gives $K\otimes_RR'\subset K'+I(K\otimes_RR')$. Nakayama on the finite quotient proves equality of kernels. For a direct-summand target, add its complementary summand to the source and apply the result already proved.
\end{proof}

\begin{lemma}\label{lem:generization}
Assume conditions \ref{ax:zariski}--\ref{ax:valuation} and $(7_{\mathrm{rel}})$. Every generization of a point of $M(T,\eta)$ belongs to $E(T,\eta)$.
\end{lemma}
\begin{proof}
Let $x\in M(T,\eta)$ map to $s\in S$. Over $\kappa(s)$ the functor is represented by finite local schemes, by the Artinian comparison. After base change, the induced pair is locally modular for the restricted functor: use the original local universal property and pair its unique map with the new base map; the minimal-field condition identifies the required point of the fibre product. Therefore $T_s$ has a local-ring isomorphism to a finite representing branch at $x$, and $T/S$ is quasi-finite there.

Put $A=\OO_{S,s}$ and pass to $\widehat A$. The branch of $T_{\widehat A}$ at a point over $x$ has a finite local open neighborhood $C=\Spec B$, by the henselian finite-branch theorem \cite[Tag 04GH]{Stacks}. Its closed point is locally modular for $F_{\widehat A}$ on $\widehat A$-schemes. Condition $(7_{\mathrm{rel}})$, applied with complete base $\widehat A$, makes it so everywhere for that same restricted functor. For any noetherian $\widehat A$-test $(Y,\zeta)$ the fibre product $h_C\times_{F_{\widehat A}}h_Y$ satisfies the sheaf, finitary and effectivity conditions. Its Artinian lifting is unique by modularity, and valuative uniqueness follows from separation of the finite scheme $C$ over the base. Lemma~\ref{lem:etale} represents it by a separated \'etale $Y$-scheme. Finitary approximation handles arbitrary affine tests, and gluing handles all tests. Hence $h_C\to F_{\widehat A}$ is representable and \'etale.

Let $X=\Spec\OO_{T,x}$. Every function invertible at $x$ is a unit on the local scheme $C$, so $C$ is an open subscheme of $X\times_A\widehat A$. The map $C\to X$ is flat local on rings and therefore fpqc. Unique formal lifting for $C$ descends to $X$ as follows. For a nilpotent immersion $Z_0\subset Z$ and a map $Z_0\to X$, the inverse image of $C$ is an open $V_0$ of $(Z_0)_{\widehat A}$. It extends uniquely to an open $V$ of $Z_{\widehat A}$, and $V\to Z$ is fpqc: flatness holds because $V$ is open in the flat base change $Z\times_A\widehat A$, while surjectivity and quasi-compactness follow from the corresponding properties after nilpotent reduction. Lift over $V$ to $C$. On the overlap, pair both resulting maps to $X$ with the base structure from the first factor. Their reductions land in the same open $C$, so the maps do as well. Formal uniqueness for $C$ makes them equal. Fpqc descent then gives the unique map over $Z$.

Thus $h_X\to F_A$ is formally \'etale. A local Artinian test centered at a generization of $x$, and any nilpotent lifting of it into $T$, factor through $X$. This gives exactly the required Artinian lifting property at that generization.
\end{proof}

\begin{lemma}\label{lem:stratum}
Assume conditions \ref{ax:zariski}--\ref{ax:valuation} and \ref{ax:generic}. If an irreducible closed subset $Z\subset T$ has generic point $z\in E(T,\eta)$, then a nonempty relatively open subset of $Z$ is contained in $E(T,\eta)$.
\end{lemma}
\begin{proof}
Over the residue field of the image $s$ of $z$, the functor is represented by finite local schemes. The infinitesimal criterion makes the fibre map \'etale at $z$, so $T/S$ is quasi-finite there. After an \'etale neighborhood of $s$ and restriction around a point over $z$, we may take $T=\Spec C$ finite over an affine base $\Spec R$. To justify the replacement, isolate the finite local branch over the strict henselization using Tag 04GH; its finite presentation and its open inclusion descend to an \'etale neighborhood of finite presentation. At the same time the finite separable residue extension of the \'etale fibre map can be split. Over a separably closed base residue field all finite local branches have purely inseparable residue fields \cite[Tag 06DD]{Stacks}, so an \'etale map between them has equal residue fields. The required finite residue splitting descends to an \'etale neighborhood. These changes have open images and preserve and reflect the lifting condition.

Let $I\subset R$ define the closure of the image of $z$, and choose $N$ from Lemma~\ref{lem:uniform} for the finite module $C$. Over the Artinian local ring $(R/I^N)_I$, the functor is represented. The branch at $z$ maps \'etale to its finite local representing branch with identical residue field; hence that local map is an isomorphism. Remove the other irreducible components from $T\times_R\Spec(R/I^N)$. On the resulting irreducible open the generic point is locally modular over the irreducible noetherian base $R/I^N$. Condition~\ref{ax:generic} supplies a nonempty open $W$ on which it is locally modular.

Fix $w\in W$ with base image $q$, and put $R'=\widehat{R_q}$. The finite algebra $C\otimes_RR'$ splits into finite local factors; let $C_w$ be the factor through $w$. Conditions \ref{ax:complete} and \ref{ax:artin} give a finite complete local $R'$-algebra $D_w$ representing the local Artinian deformation functor at the selected field point. Here finiteness of the inverse limit follows by lifting a finite set of module generators at level one and correcting their coefficients successively in powers of the maximal ideal. Completeness gives a surjection from a finite free $R'$-module, and identifies its reductions with the finite-level representing algebras. The element on $C_w$ induces a homomorphism $D_w\to C_w$.

Local modularity over $R/I^N$ identifies these two complete local deformation algebras modulo $I^N$. Lemma~\ref{lem:uniform}, including its direct-summand assertion for the target $C_w$, makes the homomorphism itself an isomorphism. Hence the original pair is Artinian formally \'etale at $w$. Indeed centered Artinian tests factor uniquely through the corresponding completions, so this conclusion is precisely the definition of $E$.

If the \'etale alternative to the conclusion of condition~\ref{ax:generic} is used, the residue extension of $D_w/I^ND_w\to C_w/I^NC_w$ is finite separable. Take the finite \'etale local extension $D'_w$ of the henselian ring $D_w$ having that residue field. Its map to $C_w$ lifts by henselian uniqueness. Modulo $I^N$ it is an isomorphism, by the unique Artinian lifting properties and completeness. Lemma~\ref{lem:uniform} makes $D'_w\to C_w$ an isomorphism. Thus $D_w\to C_w$ is \'etale, which gives the same conclusion.

For a finite morphism, primes over $I$ are incomparable. Thus near $z$ the chosen irreducible component of the inverse image of $V(I)$ is exactly the transformed stratum $Z$. The open $W$ supplies the required relative open after the modifications, and its \'etale image gives a nonempty relative open in the original $Z$.
\end{proof}


\begin{proof}[Proof of Theorem~\ref{thm:openness}]
Work on a noetherian affine piece of $T$ over an affine piece of $S$, and let $U$ be the interior of $E=E(T,\eta)$. We first prove $M(T,\eta)\subset U$. If a modular point $x$ lay outside $U$, choose an irreducible component $Z$ of the closed complement containing it, with generic point $z$. Lemma~\ref{lem:generization} puts $z$ in $E$. Lemma~\ref{lem:stratum} gives a nonempty relative open $W\subset Z$ lying in $E$. Delete the other components of the complement and shrink around $z$ until its intersection with the complement lies in $W$. This open neighborhood lies entirely in $E$: its points lie either in $U$ or in $W$. By definition of the interior it lies in $U$, a contradiction.

On $U$, each fibre product $h_U\times_Fh_Y$ is represented and \'etale by Lemma~\ref{lem:etale}, first for noetherian $Y$, then for arbitrary tests by finitarity. Its image $Q$ is an open subfunctor of $F$, and the \'etale equivalence relation $U\times_FU\rightrightarrows U$ presents an algebraic space $Q$. Every point of $U$ is quasi-finite over $S$, as is seen on the represented Artinian residue fibre. Thus $Q/S$ is locally quasi-finite.

We verify separation before using the scheme criterion. For affine chart pieces $U_i,U_j$, the relation $U_i\times_QU_j\to U_i$ is separated and \'etale, and its fibres are finite because their points inject into a fibre of the finite-type quasi-finite scheme $U_j/S$. A separated \'etale map with finite fibres over a noetherian scheme is quasi-compact. In fact the loci where its fibre degree is at least $d$ are open \cite[Tag 086R]{Stacks}; the finite degrees at the finitely many generic points bound all degrees. Tag 07RY then gives a finite locally closed stratification on which the map is finite locally free. Each inverse image is quasi-compact, proving the claim. Hence the relation is quasi-compact and $Q/S$ is quasi-separated. Valuative uniqueness makes it separated \cite[Tag 03KV]{Stacks}. A separated locally quasi-finite algebraic space over a scheme is a scheme \cite[Tag 0417, Proposition 67.50.2]{Stacks}.

On $U$, modularity is now the local-ring-isomorphism locus of $U\to Q$. This is open by finite presentation and spreading of the inverse, as in Proposition~\ref{prop:finite-type}. Since every modular point lies in $U$, this proves the assertion on the affine piece, hence globally. In the \'etale alternative to $(7_{\mathrm{rel}})$, Lemma~\ref{lem:generization} already has exactly the \'etale chart it needs; the preceding proof is unchanged.
\end{proof}

\section{The locally quasi-finite analogue}\label{sec:criterion}



The finiteness in condition~\ref{ax:artin} allows formal branches to be assembled as finite algebras. The strengthened propagation premise makes their restrictions to the punctured base \'etale. Valuative uniqueness prevents different branches from meeting there, and the openness theorem turns local representations into global charts. This explains where each hypothesis enters the proof.

\begin{lemma}\label{lem:local-global}
Suppose conditions \ref{ax:zariski}--\ref{ax:valuation}, $(7_{\mathrm{rel}})$, and \ref{ax:generic} hold and $F$ is represented by a separated locally quasi-finite scheme over each $\Spec\OO_{S,s}$. Then it is represented by such a scheme over $S$.
\end{lemma}
\begin{proof}
Take an affine finite-type open of a local representing scheme. Its finite presentation and universal element descend, by finitarity, to an affine finite-type pair $(V,\xi)$ over a neighborhood of $s$. Over the local base it is modular at every point. These points are also modular over $S$: a local $S$-test centered at a generization of $s$ factors uniquely through $\Spec\OO_{S,s}$. Theorem~\ref{thm:openness} makes its modularity locus open in $V$. Its closed complement has empty localization at $s$; finite presentation and quasi-compactness allow one denominator outside $s$ to remove it. After shrinking the base, $h_V\to F$ is modular everywhere and in particular representable \'etale, by Lemma~\ref{lem:etale} applied to its fibre products.

Taking these charts for all local points gives a surjective \'etale presentation of $F$ and hence an algebraic space. It is locally quasi-finite. The relation between two affine chart pieces has finite fibres over either piece, checked on the local representing schemes. The finite-fibre argument in Theorem~\ref{thm:openness} makes the relation quasi-compact. Valuative uniqueness gives separation, and Tag 0417 gives a scheme. This proves the lemma.
\end{proof}

\begin{proof}[Proof of Theorem~\ref{thm:criterion}]
We prove sufficiency by induction on the dimension of noetherian bases, simultaneously after their noetherian base changes. Over an Artinian base, condition~\ref{ax:artin} and Lemma~\ref{lem:comparison} give the representation. Suppose the assertion is known in dimensions less than $n$, and first take a complete noetherian local base $(A,\mathfrak m)$ of dimension $n$. Its punctured spectrum $U$ has dimension at most $n-1$, so $F_U$ is represented by a separated locally quasi-finite $X_U$.

For each component of the represented closed fibre, condition~\ref{ax:artin} gives compatible finite local algebras $B_{i,r}$ over $A/\mathfrak m^r$. Their limit $B_i$ is finite local over $A$: lift generators from the first level and correct coefficients successively in powers of $\mathfrak m$ as in Lemma~\ref{lem:stratum}. This also proves $B_i/\mathfrak m^rB_i=B_{i,r}$. Condition~\ref{ax:complete}, applied to the finite complete local ring $B_i$, supplies the universal element $\xi_i$. The two adic filtrations on this ring are cofinal. Condition $(7_{\mathrm{pro}})$ makes the maps $Z_i=\Spec B_i\to F$ \'etale, in particular over $U$.

The maps $(Z_i)_U\to X_U$ are open immersions with pairwise disjoint images. Indeed, an off-diagonal point of their equality relation has closure in the finite $A$-scheme $Z_i\times_AZ_j$ meeting the closed fibre. A noetherian valuative specialization gives a trait from a point of that relation to a point in the closed fibre, after allowing a residue extension. The two sections of $F$ agree at its generic point and hence everywhere by condition~\ref{ax:valuation}. At the closed point their Artinian universal properties force the same closed-fibre component, and for $i=j$ force the two maps $B_i\to V/\mathfrak m_V^r$ to agree for every $r$. Krull intersection in the discrete valuation ring $V$ makes their maps to $V$ equal, contradicting the chosen off-diagonal generic point. Each map $(Z_i)_U\to X_U$ is separated: its source is separated over $U$, since $Z_i$ is finite over $A$, and $X_U$ is an $U$-scheme. Thus its diagonal, already open by \'etaleness, is also closed. Hence the equality relations have no complement to the diagonal as schemes. The maps are \'etale monomorphisms, thus open immersions, and different images are disjoint.

Glue $\coprod_iZ_i$ to $X_U$ along these open immersions. The resulting scheme $X$ is locally of finite type and carries a section of $F$. Its natural map $h_X\to F$ is bijective on local Artinian tests: tests centered at the closed point use the $B_{i,r}$, and the others factor through $U$. Lemma~\ref{lem:comparison} makes it an isomorphism on all tests. It is locally quasi-finite. The gluing has quasi-compact intersections on noetherian affine pieces, since distinct finite branches are disjoint over $U$; thus it is quasi-separated. Valuative uniqueness gives separation.

For a noetherian local base that is not complete, its faithfully flat completion gives this representation and a descent datum. Murre's Appendix, Theorem, case (d) \cite[\S4, pp.~259--260]{Murre}, makes that datum effective: the cover is fpqc, its source is noetherian, and the represented scheme is separated, locally quasi-finite and locally of finite presentation. Equivalently, a separated locally quasi-finite morphism is ind-quasi-affine, and such morphisms have effective fpqc descent \cite[Tags 0AP9 and 0APK]{Stacks}. These properties descend. Thus all local representations in dimension $n$ exist. Condition $(7_{\mathrm{pro}})$ implies $(7_{\mathrm{rel}})$, so Lemma~\ref{lem:local-global} gives the global assertion and completes the induction. Every noetherian local ring has finite dimension. Applying the same local-to-global lemma therefore also treats locally noetherian bases whose dimensions are not globally bounded.

Conversely let $F=h_X$ with $X/S$ separated and locally quasi-finite. The sheaf, finitary, complete-local, Artinian and valuative conditions are the corresponding standard properties of this morphism. For $(7_{\mathrm{pro}})$, over a complete local base the branch of $X$ at a closed-fibre point is finite local and complete by the henselian finite-branch theorem. A finite algebra with the same Artinian universal property has the same quotients at every level, so completeness identifies it with that branch. Its map to $X$ is an open immersion, hence locally modular everywhere. Condition~\ref{ax:generic} follows from finite presentation and spreading a local-ring isomorphism. This proves necessity.
\end{proof}

The argument realizes Murre's announced locally quasi-finite analogue with a specified pro-modular propagation axiom. It uses finite algebras, fpqc completion descent and a uniform Artin--Rees test. It does not invoke Artin approximation over a nonexcellent base.

\section{Finite descent in every rank}\label{sec:finite-descent}

We prove the finite semilocal descent theorem by reducing matrix cocycles to additive ones through successive nilpotent quotients. The key normalization makes every correcting matrix congruent to the identity modulo the radical. It works for arbitrary direct-sum rank and requires no division by an integer.

Write $S=A'$ and $C^q=S^{\otimes_A(q+1)}$. For $h\in S$ put $d(h)=h_2-h_1$ in $C^1$. An additive cocycle $z$ satisfies $z_{23}-z_{13}+z_{12}=0$ in $C^2$; a matrix cocycle $g$ satisfies $g_{12}g_{23}=g_{13}$. A matrix coboundary has the form $h_1^{-1}h_2$. These are the additive and matrix Amitsur cohomology groups and pointed sets of the statement. None of these definitions assumes flatness of $S$ over $A$.

First observe that replacing $A$ by the equalizer $B=H^0(S/A)$ changes none of the tensor powers. Indeed, for $b\in B$ the relation $b\otimes1=1\otimes b$ already holds in $S\otimes_AS$. Thus
\begin{equation}\label{eq:equalizer-tensor}
 S\otimes_AS=S\otimes_BS,
\end{equation}
and the same argument in neighboring slots gives the assertion for all higher tensor powers. In particular the descent categories are identical. The ring $B$ is a finite $A$-module, since it is a submodule of $S$; it is noetherian and semilocal. The finite injective homomorphism $B\to S$ induces a surjection on spectra.

\subsection{The nilpotent calculation}

We establish the equivalence of the two vanishing conditions on every finite initial sequence of radical quotients. The Artinian local factors can be treated separately; a factor on which $S$ is zero contributes nothing. We may therefore take $A_n$ local, with maximal ideal denoted again by $\mathfrak m$, and with $S_n/\mathfrak mS_n$ nonzero. At level zero this is a nonzero finite algebra over the residue field $k$. It is faithfully flat, so ordinary faithfully flat descent gives additive vanishing, triviality of all free-module cocycles, and
\begin{equation}\label{eq:field-equalizer}
 H^0(S_0/k)=k.
\end{equation}

Suppose first that matrix cohomology vanishes in one fixed finite positive rank $r$ at every level. We prove additive vanishing by induction. Given an additive cocycle $z$ at level $n$, subtract the differential of a lift of its primitive at level $n-1$. Its entries now lie in $\mathfrak m^n C_n^1$. This ideal is square-zero, so
\[
 g=(1+z)I_r
\]
is a matrix cocycle. Write $g=h_1^{-1}h_2$. Modulo $\mathfrak m$, the matrix $h$ is invariant, hence is an invertible matrix over $k$ by \eqref{eq:field-equalizer}. Lift it to $\GL_r(A_n)$ and multiply $h$ on the left by its inverse. This preserves the coboundary and arranges $h=I_r+Y$ with entries of $Y$ in $\mathfrak mS_n$. Consequently
\[
 h_2-h_1=z h_1=zI_r,
\]
because $\mathfrak m^{n+1}=0$. The $(1,1)$ entry gives $z=d(h_{11})$. This proves the converse implication even when the fixed rank is one.

Now assume additive vanishing at every level. We prove matrix vanishing in every cardinal rank. For a cardinal $\Lambda$, a matrix means a column-finite matrix acting on the direct sum $S_n^{(\Lambda)}$; its invertibility means that it is an automorphism of that module. Reducing a cocycle modulo $\mathfrak m$ and applying faithfully flat descent over $k$ gives a vector space over $k$. Its dimension is $\Lambda$, as can be checked after tensoring with a residue field of $S_0$. Thus its reduced cocycle is a coboundary.

An invertible column-finite matrix modulo a nilpotent ideal lifts to an invertible matrix. Lift it and its inverse with the same finite supports in each column. The errors in their two products lie in a nilpotent ideal of the endomorphism ring, and finite geometric sums correct them. This ideal is nilpotent uniformly in $\Lambda$: a product of $h$ matrices with entries in an ideal $I$ has entries in $I^h$, since each column involves only finite sums. We can therefore normalize our cocycle as $g=I+X$ with entries in $\mathfrak m C_n^1$.

Suppose inductively that the entries of $X$ lie in $\mathfrak m^q C_n^1$, where $q\geq1$. Modulo $\mathfrak m^{q+1}$ the cocycle equation is additive. The assumed additive vanishing at this quotient supplies entrywise primitives $Y$ with
\[
 dY=X\pmod{\mathfrak m^{q+1}}.
\]
Choose a zero primitive for each zero entry of $X$. The resulting matrix is column-finite. Modulo $\mathfrak m$, each entry of $Y$ is invariant and hence lies in $k$ by \eqref{eq:field-equalizer}. Subtract lifts of these scalars from $A_n$. Their differentials vanish, so the displayed equation is preserved, and now all entries of $Y$ lie in $\mathfrak mS_n$. Thus $h=I+Y$ is invertible. The changed cocycle is $h_1gh_2^{-1}$, and
\[
 h_1g-h_2=X+Y_1-Y_2+Y_1X
             \equiv0\pmod{\mathfrak m^{q+1}}.
\]
Here $Y_1X$ has the required order because $Y$ lies in $\mathfrak m$ and $X$ lies in $\mathfrak m^q$. There is no need to choose $Y$ in $\mathfrak m^q$. Nilpotence finishes this finite iteration. It proves triviality of all cardinal-rank free-module cocycles on every Artinian quotient.

\subsection{Effectiveness over Artinian equalizers}

We explain why the matrix statement proves effectiveness, including the assertion about arbitrary flat modules. A flat module over an Artinian local ring is free in arbitrary rank. To see this, lift a basis modulo the maximal ideal to obtain a map from a free module. Nilpotent Nakayama, valid without finite generation, makes the map surjective. Flatness makes its kernel have zero reduction modulo the maximal ideal; nilpotent Nakayama then makes the kernel zero.

Let $T$ be an Artinian quotient of $A$ and put $S_T=S\otimes_AT$ and $B_T=H^0(S_T/T)$. These are Artinian rings. A flat $S_T$-module is free on each of its finitely many local factors. A descent isomorphism forces its cardinal ranks to agree on factors above the same factor of $B_T$: the tensor product of their residue fields over the residue field of that $B_T$-factor is nonzero, and restriction to any field quotient identifies the two vector-space dimensions. Hence the module has one free cardinal rank on the part above each connected component of $\Spec B_T$.

By \eqref{eq:equalizer-tensor} its datum is a matrix cocycle of that rank. The preceding calculation trivializes it, so the module descends to the free module of the same rank on the corresponding component of $B_T$. A cocycle on one component can be extended by the identity to the others, allowing the proved matrix assertion to be applied component by component. A compatible homomorphism between two such descended modules is a column-finite matrix whose entries belong to the equalizer $B_T$. This proves full faithfulness as well as essential surjectivity.

Conversely, if all flat modules descend effectively, a datum on $S_T^r$ descends to a flat $B_T$-module. On each Artinian local component it is free, and its rank is $r$, as is detected after tensoring with a residue field of $S_T$. Choosing bases makes the original cocycle a coboundary. We have proved all three equivalences in Theorem~\ref{thm:finite-descent}, including their independence of the fixed finite rank.

\subsection{The semilocal ring before completion}

We next deduce the finite-rank assertions over $A$ itself. For an additive cocycle $z\in C^1$, vanishing modulo every power of $\mathfrak m$ says
\[
 z\in\operatorname{im}(d:C^0\to C^1)+\mathfrak m^N C^1
 \quad\text{for every }N.
\]
The quotient $C^1/\operatorname{im}d$ is a finite $A$-module. Krull intersection, since $\mathfrak m\subset\operatorname{rad}A$, makes the intersection of its $\mathfrak m$-powers zero. Thus $z$ is a boundary.

For a finite-rank matrix cocycle $g$, consider the $A$-linear map between finite $A$-modules
\[
 T:M_r(S)\longrightarrow M_r(S\otimes_AS),
 \qquad T(h)=h_2-h_1g.
\]
Artin--Rees provides an integer $c$ such that
\[
 \operatorname{im}T\cap\mathfrak m^{c+1}M_r(S\otimes_AS)
                    \subset\mathfrak m\operatorname{im}T.
\]
Trivialize $g$ modulo a sufficiently large power of $\mathfrak m$ and lift the resulting matrix to $h\in M_r(S)$. Its error $T(h)$ lies in the intersection on the left. There is $w\in\mathfrak mM_r(S)$ with $T(w)=T(h)$, so $h-w$ solves the equation exactly. Its reduction modulo $\mathfrak mS$ is invertible. Since $S$ is finite over the semilocal ring $A$, the ideal $\mathfrak mS$ is contained in the radical of $S$, and $h-w$ is invertible. Hence $g$ is a coboundary. This argument does not require compatible choices of bases on all quotients, or completeness of $A$.

Finally, a finite flat $S$-module has a locally constant finite rank function. Descent identifies the ranks at any two points above the same point of $\Spec B$, by tensoring residue fields as before. This invariant rank function descends to a locally constant function on $\Spec B$: the finite surjective map is closed, so an invariant clopen set and its complement have closed images. Decompose $B$ by the finitely many resulting clopen rank sets. On each piece the module has constant finite rank and is free, since $S$ is semilocal. Matrix vanishing descends its datum to a free $B$-module. The equalizer on matrix entries proves full faithfulness. Thus $\Spec S\to\Spec B$ is effective descent for finite flat modules.

\subsection{Flat base change and arbitrary rank}

Assume now that $A$ is Artinian. Let $R$ be a flat $B$-algebra and put $S_R=S\otimes_BR$. Suppose $M$ is a flat $S_R$-module carrying a descent datum for $S_R/R$. It is flat over $S$, since $S_R$ is flat over $S$. After forgetting the $R$-action, the datum is one for $S/B$: indeed
\[
 (S\otimes_BS\otimes_BR)\otimes_{S\otimes_BR}M
       =(S\otimes_BS)\otimes_SM,
\]
with the appropriate first or second structure map, and the analogous identity holds for the triple tensor product. Artinian effectiveness gives a flat $B$-module $N$ and an isomorphism $S\otimes_BN\simeq M$ compatible with descent.

Multiplication by each $r\in R$ is a morphism of descent data. Full faithfulness descends it uniquely to an endomorphism of $N$. Uniqueness preserves products, sums and the identity, so these endomorphisms make $N$ an $R$-module. The displayed isomorphism is $R$-linear and therefore identifies $S_R\otimes_RN$ with $M$.

It remains to check flatness over $R$. Work on one local factor of $B$, write $\mathfrak b$ for its radical and $k$ for its residue field, and set $R_0=R/\mathfrak bR$ and $S_0=S/\mathfrak bS$. The nonzero finite $k$-algebra $S_0$ is faithfully flat over $k$. Since $M$ is flat over $S_R$, its quotient
\[
 M/\mathfrak bM=(S_0\otimes_kR_0)\otimes_{R_0}(N/\mathfrak bN)
\]
is flat over $S_0\otimes_kR_0$. Faithfully flat descent of flatness makes $N/\mathfrak bN$ flat over $R_0$.

Both $R$ and $N$ are flat over $B$. A free $B$-resolution of $k$, tensored with $R$, is consequently a free $R$-resolution of $R_0$, and gives
\[
 \operatorname{Tor}_i^R(R_0,N)=\operatorname{Tor}_i^B(k,N)=0
 \qquad(i>0).
\]
These facts imply $R$-flatness without any finiteness assumption on $N$ or $R$. For completeness, change of rings gives $\operatorname{Tor}_1^R(Q,N)=0$ for every $R_0$-module $Q$. Every $R$-module has a finite filtration by powers of the nilpotent ideal $\mathfrak bR$, whose quotients are $R_0$-modules. The Tor exact sequences then give the same vanishing for every $R$-module $Q$, which is the flatness criterion.

Full faithfulness over $R$ follows from that over $B$ and the unique descended $R$-actions. This proves effectiveness after every flat affine base change. Affine localization and uniqueness glue the conclusion for an arbitrary flat $B$-scheme. The proof of Theorem~\ref{thm:finite-descent} is complete.


\section{Further questions}

The strengthened propagation axiom specifies what the finite formal branches must do away from the closed point. A remaining application problem is to derive this condition from independent geometric hypotheses for particular deformation functors. Effectivity for all noetherian adic completions does not do so, by Theorem~\ref{thm:adic}. The finite descent theorem gives arbitrary-rank effectiveness over Artinian bases and their flat base changes; the proof over a general semilocal noetherian base uses finite matrices and does not give an arbitrary-rank extension there.

\paragraph{Method and availability.}

GPT-6 Astra (ChatGPT Pro mode) produced the mathematics and text. Separate GPT-6 Astra proof reconstructions and adversarial readings are recorded with their exact statements, scope and proof exposure. Earlier auxiliary-assisted checks and Claude Opus 5.5 checks at maximum effort are distinguished in those records. No human has checked the paper. Original text is CC0~1.0. Sources and checks: \url{\workbenchurl}; permanent records: \url{\recordurl}.

\enlargethispage{1.5\baselineskip}
\begin{thebibliography}{99}
\footnotesize\setlength{\itemsep}{0pt}
\bibitem{Groth131} A.~Grothendieck, Montpellier archives, cote 131, sheets 13, 34, 36 (PDF pages 14, 35, 37), Universit\'e de Montpellier. \url{https://grothendieck.umontpellier.fr/131.pdf}.
\bibitem{Groth128} A.~Grothendieck, Montpellier archives, cote 128, sheets 25--29, Universit\'e de Montpellier. \url{https://grothendieck.umontpellier.fr/128.pdf}.
\bibitem{Murre} J.~P.~Murre, Representation of unramified functors. Applications (according to unpublished results of A.~Grothendieck), \emph{S\'eminaire Bourbaki}, expos\'e 294, May 1965, in vol.~9 (1964--1966), 243--261. \url{https://www.numdam.org/item/SB_1964-1966__9__243_0/}.
\bibitem{dJO} A.~J.~de Jong and N.~Olander, On weakly \'etale morphisms, \emph{J. Pure Appl. Algebra} \textbf{228} (2024), 107532. Theorem 2, p.~2. \url{https://doi.org/10.1016/j.jpaa.2023.107532}.
\bibitem{FeldPurity} N.~Feld, \emph{Gersten injectivity and Morel's purity conjecture}, AI-assisted preprint, October--November 2026. Theorem~B, p.~2; AI disclosure, \S9, p.~15. \url{https://hexagonmath.org/2611.00017v1}.
\bibitem{Stacks} The Stacks Project Authors, \emph{The Stacks Project}, 2026. \url{https://stacks.math.columbia.edu}.
\end{thebibliography}
\end{document}
